% Architecture: Visual Guide to the TIDAL Lagrangian-to-PDE Pipeline
% Source: docs/architecture/README.md (migrated to LaTeX)
% Uses macros from preamble.tex

\section{Architecture}
\label{sec:architecture}

Visual guide to the TIDAL Lagrangian-to-PDE pipeline.
All PDEs derive from a Lagrangian via symbolic computation --- no physics is hardcoded in Python.

\subsection{Data Flow Pipeline}
\label{sec:data-flow-pipeline}

The pipeline has two stages separated by a JSON specification file.

\textbf{Stage~1 (Mathematica/xAct)} derives field equations symbolically from a Lagrangian.
\textbf{Stage~2 (Python/SUNDIALS + numpy)} loads the specification and runs a numerical simulation.

\begin{figure}[htbp]
\centering
\documentclass[border=10pt]{standalone}
\usepackage{tidal-tikz-styles}

\begin{document}
\begin{tikzpicture}[
    node distance=0.5cm
]

% ===================================================================
% ROW 1: Main pipeline stages (left to right)
% ===================================================================

% --- Input ---
\node[data] (toml) {\texttt{theory.toml}};

% --- Derive ---
\node[stage=Pink, right=1.0cm of toml] (derive)
    {\texttt{tidal derive}\\[2pt]\scriptsize Wolfram/xAct};

% --- JSON (narrow waist) ---
\node[data, right=1.0cm of derive, minimum width=1.8cm,
      draw=BlueLine, fill=BlueFill, fill opacity=0.35,
      font=\footnotesize\bfseries] (json)
    {\texttt{spec.json}};

% --- Simulate ---
\node[stage=Purple, right=1.0cm of json] (sim)
    {\texttt{tidal simulate}\\[2pt]\scriptsize PDE solver};

% --- Snapshots ---
\node[data, right=1.0cm of sim] (snaps)
    {snapshots\\[-1pt]\scriptsize\texttt{.npy}};

% --- Measure ---
\node[stage=Green, right=1.0cm of snaps] (meas)
    {\texttt{tidal measure}\\[2pt]\scriptsize 13 metrics};

% --- Results ---
\node[data, right=1.0cm of meas] (results)
    {metrics\\[-1pt]\scriptsize\texttt{.csv}};

% ===================================================================
% ARROWS (main flow)
% ===================================================================
\draw[arr] (toml) -- (derive);
\draw[arr] (derive) -- (json);
\draw[arr] (json) -- (sim);
\draw[arr] (sim) -- (snaps);
\draw[arr] (snaps) -- (meas);
\draw[arr] (meas) -- (results);

% ===================================================================
% ROW 2: External libraries (below stages)
% ===================================================================
\node[lib, below=0.15cm of derive] (lib1)
    {xAct, xPert, xCoba};
\node[lib, below=0.15cm of sim] (lib2)
    {SUNDIALS, NumPy\\SciPy};
\node[lib, below=0.15cm of meas] (lib3)
    {NumPy, SciPy\\FFT};

% ===================================================================
% ROW 3: Validation checks (below libraries)
% ===================================================================
\node[check, below=0.55cm of derive] (chk1)
    {\texttt{\_validate\_config()}\\[1pt]
     \texttt{\_validate\_lagrangian()}};

\node[check, below=0.55cm of json] (chk2)
    {\texttt{from\_dict()}\\[1pt]
     field/operator/matrix\\validation};

\node[check, below=0.55cm of sim] (chk3)
    {\texttt{ensure\_consistent\_ic()}\\[1pt]
     constraint pre-solve};

% Check arrows (dotted, pointing up)
\draw[dataarr, densely dotted] (chk1) -- (derive);
\draw[dataarr, densely dotted] (chk2) -- (json);
\draw[dataarr, densely dotted] (chk3) -- (sim);

% ===================================================================
% TOP: Sweep/Plot orchestration (above main flow)
% ===================================================================
\node[stage=Lime, above=1.2cm of sim, minimum width=2.0cm,
      minimum height=1.0cm] (sweep)
    {\texttt{tidal sweep}\\[2pt]\scriptsize parameter grid};

% Sweep arrows: wraps around simulate+measure
\draw[arr, rounded corners=4pt]
    (sweep.west) -| ([xshift=-0.3cm]sim.north) -- (sim.north);
\draw[arr, rounded corners=4pt]
    (meas.north) -- ([yshift=0.3cm]meas.north) -| (sweep.east);

% Plot branch
\node[stage=Lime, right=0.8cm of results, minimum width=1.6cm,
      minimum height=1.0cm] (plot)
    {\texttt{tidal plot}};
\draw[arr] (results) -- (plot);

% ===================================================================
% BRACE: Symbolic vs Numerical separation at JSON
% ===================================================================
\draw[decorate, decoration={brace, amplitude=6pt, raise=2pt},
      PinkLine, thick]
    (toml.north west) -- (derive.north east)
    node[midway, above=10pt, font=\scriptsize\bfseries, text=PinkLine]
    {Symbolic (CAS)};

\draw[decorate, decoration={brace, amplitude=6pt, raise=2pt},
      PurpleLine, thick]
    (sim.north west) -- (meas.north east)
    node[midway, above=10pt, font=\scriptsize\bfseries, text=PurpleLine]
    {Numerical (PDE)};

% JSON label
\node[above=0.25cm of json, font=\scriptsize\bfseries, text=BlueLine]
    {portable contract};

\end{tikzpicture}
\end{document}

\caption{TIDAL data flow pipeline. Stage~1 (Mathematica/xAct) derives field equations symbolically from the Lagrangian; the JSON specification bridges the two stages; Stage~2 (Python/SUNDIALS + numpy) loads the spec and runs numerical simulation.}
\label{fig:pipeline-overview}
\end{figure}

\subsubsection{Wolfram Module Roles}
\label{sec:wolfram-module-roles}

\begin{table}[htbp]
\centering
\caption{Wolfram pipeline modules and their roles.}
\label{tab:wolfram-modules}
\begin{tabular}{@{}lll@{}}
\toprule
Module & Purpose & Key Export \\
\midrule
\texttt{CommonUtilities.wl} & Shared helpers: CD$\to$Derivative conversion, Christoffel/epsilon evaluation, metric components & \texttt{ConvertCDToDerivatives}, \texttt{EvaluateChristoffelComponents} \\
\texttt{EulerLagrange.wl} & Derive equations of motion from a Lagrangian & \texttt{EulerLagrangeEquation[L, field, cd]} \\
\texttt{ComponentDecompose.wl} & Convert tensor EOM to scalar component equations & \texttt{DecomposeToComponents[eom, field, chart, additionalFields]} \\
\texttt{ExportJSON.wl} & Serialize component equations to JSON & \texttt{BuildMultiFieldJSONStructure[fieldEquations, metadata]} \\
\texttt{GaugeFix.wl} & Gauge fixing presets and custom expressions & \texttt{AddGaugeFixingTerm}, \texttt{BuildLorenzGaugeTerm}, \texttt{BuildDeDonderGaugeTerm} \\
\texttt{Linearize.wl} & Perturbation theory (xPert) for linearized gravity & \texttt{LinearizeTensorExpression[expr]} \\
\bottomrule
\end{tabular}
\end{table}

\subsubsection{JSON Specification Structure}
\label{sec:json-spec-structure}

\begin{lstlisting}[style=json]
{
  "spacetime": { "dimension": 3, "signature": [-1, 1, 1],
                 "coordinates": ["t", "x", "y"] },
  "fields":    [{ "name": "phi_0", "slot_kind": "field",
                   "time_derivative_order": 2 }],
  "equations": [{
    "field_name": "phi_0",
    "lhs": { "expression": "...", "order": { "time": 2 } },
    "rhs": [{ "coefficient": -1.0,
              "coefficient_symbolic": "-m2",
              "operator": "identity", "field": "phi_0",
              "time_dependent": false,
              "coordinate_dependent": false }]
  }],
  "hamiltonian_terms": [{ "coefficient": 0.5,
    "operator": "identity",
    "fields": ["v_phi_0", "v_phi_0"] }],
  "coupling":  { "mass_matrix": [[1.0]],
                 "coupling_matrix": [[0.0]] },
  "metadata":  { "lagrangian": "...", "gauge": null,
                 "linearized": false }
}
\end{lstlisting}

\subsection{Python Module Dependencies}
\label{sec:python-module-deps}

The \texttt{symbolic/} subpackage loads JSON specs into \texttt{EquationSystem} dataclasses.
The \texttt{solver/} subpackage provides the numerical engine: state management, spatial operators, RHS evaluation, and four time-integration backends.
The \texttt{cli/} subpackage provides the \texttt{tidal} command-line interface with 9 subcommands (derive, simulate, inspect, measure, list, validate, plot, sweep, analyze).
The \texttt{measurement/} subpackage provides 13 post-hoc analysis types (energy, conversion, mixing, spectrum, spectral\_conversion, dispersion, conservation, effective\_mass, asymptotic, peak\_conversion, velocity, resonance, summary).

\paragraph{Symbolic subpackage.}
\label{par:symbolic}
\begin{itemize}
  \item \texttt{json\_loader.py} --- \emph{EquationSystem, OperatorTerm, ComponentEquation, LHSStructure}
  \item \texttt{reduction.py} --- \emph{Plane-wave dimensional reduction}
\end{itemize}

\paragraph{Solver subpackage.}
\label{par:solver}
\begin{itemize}
  \item \texttt{ida.py} --- \emph{SUNDIALS IDA (DAE)}
  \item \texttt{cvode.py} --- \emph{SUNDIALS CVODE (BDF)}
  \item \texttt{leapfrog.py} --- \emph{St\"ormer-Verlet symplectic}
  \item \texttt{scipy\_solver.py} --- \emph{solve\_ivp wrapper}
  \item \texttt{fields.py} --- \emph{FieldSet (typed container)}
  \item \texttt{state.py} --- \emph{StateLayout (frozen dataclass)}
  \item \texttt{grid.py} --- \emph{GridInfo}
  \item \texttt{operators.py} --- \emph{numpy spatial operators}
  \item \texttt{coefficients.py} --- \emph{CoefficientEvaluator (4-level cache)}
  \item \texttt{rhs.py} --- \emph{RHSEvaluator}
  \item \texttt{constraint\_solve.py} --- \emph{Three-tier constraint pre-solve}
\end{itemize}

\paragraph{CLI subpackage.}
\label{par:cli}
\begin{itemize}
  \item \texttt{\_\_init\_\_.py} --- \emph{9 subcommands}
  \item \texttt{\_derive.py} --- \emph{TOML $\to$ .wls $\to$ JSON}
  \item \texttt{\_simulate.py} --- \emph{JSON $\to$ solver $\to$ output}
  \item \texttt{\_measure.py} --- \emph{Measurement dispatcher}
  \item \texttt{\_sweep.py} --- \emph{Parameter sweeps}
  \item \texttt{\_analyze.py} --- \emph{Sensitivity analysis}
\end{itemize}

\paragraph{Measurement subpackage.}
\label{par:measurement}
\begin{itemize}
  \item \texttt{\_energy.py} --- \emph{Hamiltonian energy}
  \item \texttt{\_conversion.py} --- \emph{Conversion $P(t)$}
  \item \texttt{\_spectral.py} --- \emph{Power spectrum}
  \item \texttt{\_diagnostics.py} --- \emph{Conservation checks}
  \item \texttt{\_velocity.py} --- \emph{Group/phase velocity}
  \item \texttt{\_sensitivity.py} --- \emph{Sobol/Morris}
\end{itemize}

\paragraph{External libraries.}
\label{par:external}
\begin{itemize}
  \item SUNDIALS~\cite{hindmarsh2005sundials} --- IDA, CVODE (via scikit-sundae)
  \item numpy / scipy
\end{itemize}

\subsection{State Structure (StateLayout + FieldSet)}
\label{sec:state-structure}

The simulation state is managed by two classes:

\begin{itemize}
  \item \textbf{\texttt{StateLayout}} (frozen dataclass in \texttt{state.py}): maps field names to slot indices, slot kinds, and time derivative orders. Computed once from the JSON spec.
  \item \textbf{\texttt{FieldSet}} (typed container in \texttt{fields.py}): a contiguous flat numpy array with zero-copy named views for each field.
\end{itemize}

The layout depends on each component's \textbf{time derivative order} in the JSON specification:

\begin{table}[htbp]
\centering
\caption{State slot allocation by time derivative order.}
\label{tab:state-slots}
\begin{tabular}{@{}llll@{}}
\toprule
Time Order & Slots Used & Slot Kinds & Rate Computation \\
\midrule
2 (wave-like) & field + velocity & \texttt{"field"} + \texttt{"velocity"} & $\mathrm{d}/\mathrm{d}t\;\text{field} = \text{velocity}$; $\mathrm{d}/\mathrm{d}t\;\text{velocity} = \text{RHS}$ \\
1 (first-order) & field only & \texttt{"field"} & $\mathrm{d}/\mathrm{d}t\;\text{field} = \text{RHS}$ \\
0 (constraint) & field only & \texttt{"constraint"} & Algebraic: solved by \texttt{constraint\_solve.py} or IDA \\
\bottomrule
\end{tabular}
\end{table}

\subsubsection{Example: Chern--Simons ($A_0$ Constraint, $A_1$ and $A_2$ First-Order)}
\label{sec:state-example-cs}

\begin{figure}[htbp]
\centering
\documentclass[border=10pt]{standalone}
\usepackage[T1]{fontenc}
\usepackage{amsmath}
\usepackage{tikz}
\usepackage[svgnames]{xcolor}
\usetikzlibrary{shapes.geometric, arrows.meta, positioning, calc, fit,
                 decorations.pathreplacing}

% ---------------------------------------------------------------
% Colour palette (matches solver_flowchart.tex)
% ---------------------------------------------------------------
\colorlet{OrangeFill}{LightSalmon}
\colorlet{OrangeLine}{DarkSalmon!30!Chocolate}
\colorlet{GreenFill}{MediumAquamarine}
\colorlet{GreenLine}{MediumAquamarine!30!Teal}
\colorlet{PurpleFill}{MediumPurple}
\colorlet{PurpleLine}{MediumPurple!30!SlateBlue}
\colorlet{PinkFill}{Thistle}
\colorlet{PinkLine}{Thistle!30!Plum}
\colorlet{BlueFill}{LightSkyBlue}
\colorlet{BlueLine}{SteelBlue}
\colorlet{GrayFill}{LightGray!60!OldLace}
\colorlet{GrayLine}{RosyBrown!50!Silver}

\begin{document}
\begin{tikzpicture}[
    slot/.style={
        rectangle, minimum width=1.4cm, minimum height=0.7cm,
        draw=black!60, line width=0.6pt, font=\footnotesize,
        align=center, text opacity=1, text=black
    },
    fieldslot/.style={slot, fill=BlueFill, fill opacity=0.4},
    velslot/.style={slot, fill=PurpleFill, fill opacity=0.3},
    conslot/.style={slot, fill=OrangeFill, fill opacity=0.4, draw=OrangeLine, dashed},
    lbl/.style={font=\footnotesize\bfseries, text=black},
    annot/.style={font=\scriptsize\itshape, text=GrayLine, align=center},
    arr/.style={-{Stealth[length=4pt, width=3pt]}, semithick, black!70},
    node distance=0cm
]

% ===================================================================
% PART A: Field generation from inputs (top)
% ===================================================================
\node[font=\normalsize\bfseries, anchor=west] (titleA) at (0, 0)
    {A. Field count from inputs};

\node[below=0.4cm of titleA.west, anchor=west,
      rectangle, rounded corners=2pt, draw=GrayLine,
      fill=GrayFill, fill opacity=0.3, inner sep=6pt,
      font=\scriptsize, align=left, text opacity=1] (table) {%
\begin{tabular}{@{}l c c c c@{}}
    \textbf{Type} & \textbf{Rank} & \textbf{Symmetry}
    & \textbf{2+1D} & \textbf{3+1D} \\[2pt]
    Scalar & 0 & --- & 1 & 1 \\
    Vector & 1 & --- & 3 & 4 \\
    Sym.\ tensor & 2 & $T_{ab}=T_{ba}$ & 6 & 10 \\
    Antisym.\ tensor & 2 & $T_{ab}=-T_{ba}$ & 3 & 6 \\
    General tensor & 2 & none & 9 & 16 \\
\end{tabular}};

% Formula annotation
\node[annot, right=0.8cm of table.north east, anchor=north west,
      text width=4.2cm] (formula) {%
    Symmetric: $\frac{(d{+}1)(d{+}2)}{2}$\\[2pt]
    Antisymmetric: $\frac{d(d{+}1)}{2}$\\[4pt]
    Gauge fixing further reduces:\\
    e.g.\ TT gauge: $10 \to 2$ (3+1D GW)\\[4pt]
    Metric signature $\eta^{00}{=}{-}1$:\\
    temporal components $\to$ constraint\\
    (no $\partial^2_t$, time\_order${}=0$)};

% ===================================================================
% PART B: State vector examples (bottom)
% ===================================================================
\node[font=\normalsize\bfseries, anchor=west] (titleB)
    at (0, -5.0) {B. State vector layout};

% --- Example 1: Coupled scalars ---
\node[lbl, below=0.5cm of titleB.west, anchor=west] (ex1lbl)
    {Coupled scalars (1+1D, 2 fields):};

\node[fieldslot, below=0.2cm of ex1lbl.west, anchor=north west] (s1a) {$\varphi$};
\node[velslot, right=of s1a] (s1b) {$v_\varphi$};
\node[fieldslot, right=of s1b] (s1c) {$\chi$};
\node[velslot, right=of s1c] (s1d) {$v_\chi$};

% N annotation
\draw[decorate, decoration={brace, amplitude=3pt, mirror, raise=2pt},
      GrayLine]
    (s1a.south west) -- (s1a.south east)
    node[midway, below=5pt, font=\scriptsize, text=GrayLine] {$N$};

% Slot indices
\node[font=\scriptsize, text=GrayLine, above=1pt of s1a] {slot 0};
\node[font=\scriptsize, text=GrayLine, above=1pt of s1b] {slot 1};
\node[font=\scriptsize, text=GrayLine, above=1pt of s1c] {slot 2};
\node[font=\scriptsize, text=GrayLine, above=1pt of s1d] {slot 3};

% --- Example 2: Proca 2+1D ---
\node[lbl, below=1.0cm of s1a.south west, anchor=west] (ex2lbl)
    {Proca (2+1D, 1 vector field):};

\node[conslot, below=0.2cm of ex2lbl.west, anchor=north west] (s2a) {$A_0$};
\node[fieldslot, right=of s2a] (s2b) {$A_1$};
\node[velslot, right=of s2b] (s2c) {$v_{A_1}$};
\node[fieldslot, right=of s2c] (s2d) {$A_2$};
\node[velslot, right=of s2d] (s2e) {$v_{A_2}$};

% Constraint annotation
\node[annot, below=0.25cm of s2a] {constraint\\(no velocity)};

\node[font=\scriptsize, text=GrayLine, above=1pt of s2a] {slot 0};
\node[font=\scriptsize, text=GrayLine, above=1pt of s2c] {slot 2};
\node[font=\scriptsize, text=GrayLine, above=1pt of s2e] {slot 4};

% --- Example 3: Gertsenshtein 1+1D ---
\node[lbl, below=1.3cm of s2a.south west, anchor=west] (ex3lbl)
    {Gertsenshtein (1+1D, 6 surviving fields):};

\node[fieldslot, below=0.2cm of ex3lbl.west, anchor=north west] (s3a) {$h_5$};
\node[velslot, right=of s3a] (s3b) {$v_{h_5}$};
\node[fieldslot, right=of s3b] (s3c) {$h_7$};
\node[velslot, right=of s3c] (s3d) {$v_{h_7}$};
\node[fieldslot, right=of s3d] (s3e) {$a_1$};
\node[velslot, right=of s3e] (s3f) {$v_{a_1}$};

\node[annot, right=0.4cm of s3f.east, anchor=west, text width=3.0cm]
    {from $h_{\mu\nu}$ (10 components)\\
     + $a_\mu$ (2 components)\\
     $\xrightarrow{\text{TT + Lorenz}}$ 6 surviving\\
     $\xrightarrow{\text{plane-wave}}$ 6 (1+1D)};

% ===================================================================
% LEGEND
% ===================================================================
\node[below=1.0cm of s3a.south west, anchor=north west,
      draw=GrayLine, rounded corners=2pt, inner sep=6pt,
      fill=white, font=\scriptsize] (legend) {%
\begin{tabular}{@{}l l@{}}
    \textbf{Slot types} \\[3pt]
    \tikz\fill[BlueFill, fill opacity=0.4, draw=black!60]
        (0,0) rectangle (0.5,0.25); & Field ($q$): time\_order $\geq 1$ \\[2pt]
    \tikz\fill[PurpleFill, fill opacity=0.3, draw=black!60]
        (0,0) rectangle (0.5,0.25); & Velocity ($v = dq/dt$): time\_order $\geq 2$ \\[2pt]
    \tikz\fill[OrangeFill, fill opacity=0.4, draw=OrangeLine, dashed]
        (0,0) rectangle (0.5,0.25); & Constraint: time\_order $= 0$ (no velocity slot) \\
\end{tabular}};

\end{tikzpicture}
\end{document}

\caption{State layout for the Chern--Simons example: $A_0$ is a constraint (solved algebraically), while $A_1$ and $A_2$ are first-order fields with rates computed by \texttt{RHSEvaluator}.}
\label{fig:state-layout}
\end{figure}

The layout maps field names to contiguous slots:
\begin{itemize}
  \item $A_0$: \texttt{order=0} $\to$ constraint --- 1 slot, solved algebraically
  \item $A_1$: \texttt{order=1} $\to$ first-order --- 1 slot, rate $= \text{RHS}$
  \item $A_2$: \texttt{order=1} $\to$ first-order --- 1 slot, rate $= \text{RHS}$
\end{itemize}

The \texttt{RHSEvaluator} output:
\begin{itemize}
  \item $A_0$: residual $= 0$ (IDA algebraic)
  \item $\mathrm{d}/\mathrm{d}t\;A_1 = \text{RHS}_1$
  \item $\mathrm{d}/\mathrm{d}t\;A_2 = \text{RHS}_2$
\end{itemize}

\subsubsection{Example: Coupled Scalars (Two Second-Order Fields)}
\label{sec:state-example-coupled}

For two second-order fields, the state vector contains four slots:
\begin{itemize}
  \item Slot~0: $\phi_0$ (field)
  \item Slot~1: $v_{\phi_0}$ (velocity)
  \item Slot~2: $\chi_0$ (field)
  \item Slot~3: $v_{\chi_0}$ (velocity)
\end{itemize}

The \texttt{RHSEvaluator} output:
\begin{itemize}
  \item $\mathrm{d}/\mathrm{d}t\;\phi_0 = v_{\phi_0}$
  \item $\mathrm{d}/\mathrm{d}t\;v_{\phi_0} = \sum \text{coeff} \times \text{op}(\phi, \chi)$
  \item $\mathrm{d}/\mathrm{d}t\;\chi_0 = v_{\chi_0}$
  \item $\mathrm{d}/\mathrm{d}t\;v_{\chi_0} = \sum \text{coeff} \times \text{op}(\phi, \chi)$
\end{itemize}

Cross-field coupling: a term in the $\phi$ equation can reference $\chi_0$ via \texttt{\{"operator": "laplacian", "field": "chi\_0"\}}.

Velocity naming convention: \texttt{v\_\{field\_name\}} (e.g., \texttt{v\_phi\_0}, \texttt{v\_A\_1}) --- Euler--Lagrange velocity form, not canonical momenta.

\subsection{Execution Flow: \texttt{RHSEvaluator.evaluate()}}
\label{sec:rhs-evaluator}

Called by the solver backend at every timestep. Iterates over field equations and accumulates RHS terms.

\begin{figure}[htbp]
\centering
\documentclass[border=10pt]{standalone}
\usepackage{tidal-tikz-styles}

\begin{document}
\begin{tikzpicture}[
    node distance=0.5cm and 1.8cm,
    thick
]

% ===================================================================
% NODES
% ===================================================================

% --- Start ---
\node[solver] (start) {%
    \textbf{Auto-select solver}\\[1pt]
    scheme, spec, grid, bc};

% --- D0: Explicit scheme? ---
\node[decision, below=of start] (d0) {Explicit\\scheme?};

% --- Explicit result ---
\node[solver, right=of d0] (explicit)
    {Use specified\\solver};

% --- D1: Modal eligible? ---
\node[decision, below=of d0] (d1) {Modal\\eligible?};

% --- Modal solver ---
\node[solver=Purple, right=of d1] (modal)
    {Modal solver};

% --- D2: Constraint eqs? ---
\node[decision, below=of d1] (d2) {Constraint\\equations?};
\node[annot, right=-0.7cm of d2.east, anchor=north west]
    (a2) {zeroth-order in $t$\\(elliptic / algebraic DAE)};

% --- IDA (constraint) ---
\node[solver=Orange, right=of d2] (ida1) {IDA};

% --- D3: First-order eqs? ---
\node[decision, below=of d2] (d3) {First-order\\equations?};
\node[annot, right=-0.4cm of d3.east, anchor=north west]
    (a3) {first-order in $t$\\(diffusion / transport) \\ (implicit BDF)};

% --- IDA (first-order) ---
\node[solver=Orange, right=of d3] (ida2) {IDA};

% --- D4: Dissipation? ---
\node[decision, below=of d3] (d4) {Dissipation?};
\node[annot, right=0.0cm of d4.east, anchor=north west]
    (a4) {dissipative $\partial_t$ terms\\(breaks symplecticity)};

% --- IDA (dissipation) ---
\node[solver=Orange, right=2.2cm of d4] (ida3) {IDA};

% --- D5: Missing Hamiltonian? ---
\node[decision, below=of d4] (d5) {Missing\\Hamiltonian?};
\node[annot, below right=0.1cm and -0.4cm of d5.east, anchor=north west]
    (a5) {no Hamiltonian available\\for second-order wave eqs};

% --- Warning + fallthrough ---
\node[solver=Green, right=1.86cm of d5] (warn)
    {\textit{Warning} $\to$\\CVODE fallback};

% --- Default: CVODE ---
\node[solver=Green, left=of d5] (cvode)
    {CVODE};
\node[annot, below left=0.1cm and -0.5cm of d5.west, anchor=north east]
    (a6) {Pure wave, Hamiltonian\\(adaptive BDF)};

% --- Modal sub-check panel (below the flowchart; references D1) ---
\node[subcheck, below=0.8cm of cvode.south west, anchor=north west] (subchecks) {%
    \textbf{Modal eligibility (D1) requires all of:}\\[3pt]
    1.\ Flat metric\\
    2.\ All BCs periodic\\
    3.\ No time-dependent coefficients\\
    4.\ All operators have exact Fourier multipliers\\
    5.\ Constraints Fourier-eliminable (when present)};

% ===================================================================
% EDGES
% ===================================================================

% Start -> D0
\draw[arr] (start) -- (d0);

% D0 -> Explicit (yes)
\draw[arr] (d0) -- node[yeslbl, above] {Yes} (explicit);

% D0 -> D1 (no / "auto")
\draw[arr] (d0) -- node[nolbl, left] {No (auto)} (d1);

% Sub-check panel -> D1 (now sits below the flowchart; reference link only)
\draw[dotarr] (subchecks.north) to[out=90, in=180]
    node[yeslbl, midway, sloped, above, text=GrayLine]
    {\itshape requires all} (d1.west);

% D1 -> Modal (yes)
\draw[arr] (d1) -- node[yeslbl, above] {Yes} (modal);

% D1 -> D2 (no)
\draw[arr] (d1) -- node[nolbl, left] {No} (d2);

% D2 -> IDA (yes)
\draw[arr] (d2) -- node[yeslbl, above] {Yes} (ida1);

% D2 -> D3 (no)
\draw[arr] (d2) -- node[nolbl, left] {No} (d3);

% D3 -> IDA (yes)
\draw[arr] (d3) -- node[yeslbl, above] {Yes} (ida2);

% D3 -> D4 (no)
\draw[arr] (d3) -- node[nolbl, left] {No} (d4);

% D4 -> IDA (yes)
\draw[arr] (d4) -- node[yeslbl, above] {Yes} (ida3);

% D4 -> D5 (no)
\draw[arr] (d4) -- node[nolbl, left] {No} (d5);

% D5 -> warn (yes)
\draw[arr] (d5) -- node[yeslbl, above] {Yes} (warn);

% D5 -> CVODE (no)
\draw[arr] (d5) -- node[nolbl, above] {No} (cvode);

\end{tikzpicture}
\end{document}

\caption{Execution flow of \texttt{RHSEvaluator.evaluate()}: for each field equation, resolve the target field, apply the spatial operator, resolve the coefficient via \texttt{CoefficientEvaluator}, and accumulate.}
\label{fig:solver-flowchart}
\end{figure}

The evaluation loop for each field equation:
\begin{enumerate}
  \item Initialize accumulator $= 0$.
  \item For each RHS term:
  \begin{enumerate}
    \item Resolve target field from \texttt{FieldSet} (zero-copy view).
    \item Apply spatial operator (numpy: laplacian, gradient, etc.).
    \item Resolve coefficient via \texttt{CoefficientEvaluator}.
    \item $\text{accumulator} \mathrel{+}= \text{coefficient} \times \text{operated}$.
  \end{enumerate}
  \item Return numpy array of RHS values.
\end{enumerate}

\subsubsection{Solver Backends}
\label{sec:solver-backends}

Each backend calls \texttt{RHSEvaluator.evaluate()} with different integration strategies:

\begin{table}[htbp]
\centering
\caption{Solver backends and their integration methods.}
\label{tab:solver-backends}
\begin{tabular}{@{}lll@{}}
\toprule
Backend & Entry Point & Integration Method \\
\midrule
\textbf{IDA} & \texttt{ida.py} & Implicit Newton iteration for DAE systems. Residual form: $F(t, y, y') = 0$. Supports sparse Jacobian and analytical Jacobian (precomputed $\mathrm{d}F/\mathrm{d}y + \mathrm{d}F/\mathrm{d}y'$). \\
\textbf{CVODE} & \texttt{cvode.py} & BDF adaptive ODE. Tolerance-controlled (\texttt{--rtol}/\texttt{--atol}). Automatic step size selection. \\
\textbf{Leapfrog} & \texttt{leapfrog.py} & St\"ormer-Verlet symplectic. KDK splitting with force caching (halves force evaluations). \\
\textbf{scipy} & \texttt{scipy\_solver.py} & \texttt{solve\_ivp} wrapper with DOP853/Radau/BDF methods. \\
\bottomrule
\end{tabular}
\end{table}

Automatic solver selection: systems with algebraic constraints $\to$ IDA; pure wave equations $\to$ CVODE or leapfrog.

\subsubsection{Analytical Jacobian (\texttt{analytical\_jacobian.py})}
\label{sec:analytical-jacobian}

For time-independent linear systems, the Jacobian $J = \mathrm{d}F/\mathrm{d}y + c_j \cdot \mathrm{d}F/\mathrm{d}y'$
is precomputed once and supplied analytically via three delivery tiers:

\begin{table}[htbp]
\centering
\caption{Analytical Jacobian delivery tiers.}
\label{tab:jacobian-tiers}
\begin{tabular}{@{}llll@{}}
\toprule
Tier & Size Range & Method & Speedup vs FD \\
\midrule
Dense & $N \leq 2{,}000$ & 2D \texttt{jacfn} callback & $5.3\times$ \\
Sparse & $2\text{K} < N \leq 200\text{K}$ & 1D CSC \texttt{jacfn} + SuperLU\_MT & IDA: $2.5\times$, CVODE: $1.2$--$1.4\times$ \\
GMRES & $N > 200\text{K}$ & \texttt{jactimes} mat-vec product & Eliminates $O(n_\text{colors})$ evals \\
\bottomrule
\end{tabular}
\end{table}

Construction uses COO triple accumulation (\texttt{\_COOAccumulator}) with a single
CSC conversion --- $O(\text{nnz})$ total instead of $O(N^2)$ LIL block-assignment.
Operator matrices use $O(\text{nnz})$ circulant construction for periodic BCs.
Falls through to FD tiers for time-dependent systems. Requires sksundae $\geq 1.1.2$.

\subsubsection{Coefficient Resolution Hierarchy (4-Level Cache)}
\label{sec:coefficient-cache}

Each RHS term has a coefficient that may be constant, time-dependent, or position-dependent.
\texttt{CoefficientEvaluator} (in \texttt{coefficients.py}) resolves coefficients via a fast-path hierarchy:

\begin{description}
  \item[L0: Pre-resolved constants] Computed once at \texttt{\_\_init\_\_}, stored as float.
  \item[L1: Expression string cache] Mathematica$\to$Python conversion, computed once.
  \item[L2: Spatial grid arrays] Position-dependent but time-independent, computed once per grid.
  \item[L3: Per-timestep dedup] Time+position dependent, deduplicated within each timestep.
  \begin{itemize}
    \item Constant: parameter lookup $\to$ float.
    \item Time-only: \texttt{eval(expr, \{t=t, params\})} $\to$ float.
    \item Position-dependent: \texttt{eval(expr, \{t, x, y, z\})} $\to$ ndarray (grid-shaped).
  \end{itemize}
\end{description}

\subsection{Working Examples (20 Total)}
\label{sec:working-examples}

\begin{table}[htbp]
\centering
\caption{Complete list of TIDAL example implementations.}
\label{tab:examples}
\begin{tabular}{@{}llll@{}}
\toprule
Example & Dim & Fields & Key Features \\
\midrule
\multicolumn{4}{@{}l}{\emph{Scalar and vector (flat spacetime)}} \\
\texttt{coupled\_scalars/} & 1+1D & $\varphi$, $\chi$ & Cross-field gradient coupling, Rabi conversion \\
\texttt{coupled\_scattering/} & 1+1D & $\varphi$, $\chi$ & Position-dependent Gaussian coupling, wave scattering \\
\texttt{scalar\_potential\_well/} & 1+1D & $\varphi$ & Background potential well, bound states \\
\texttt{chern\_simons/} & 2+1D & $A_\mu$ & Maxwell + topological term, $\varepsilon$ tensor, constraint \\
\texttt{elasticity/} & 2+1D & $u_x$, $u_y$ & Navier--Cauchy, anisotropic Laplacian, cross derivatives \\
\texttt{coupled\_proca/} & 2+1D & $A_\mu$, $B_\mu$ & Two massive vectors, coupled Helmholtz constraints \\
\texttt{scalar\_vector\_coupling/} & 2+1D & $\varphi$, $A_\mu$ & Mixed-rank cross-field (scalar+vector), 4 constants \\
\texttt{proca\_background/} & 2+1D & $A_\mu$, $B_\mu$ & Lorentzian background, constraint+BG integration \\
\midrule
\multicolumn{4}{@{}l}{\emph{Curved spacetime}} \\
\texttt{sphere\_kg/} & 2+1D & $\varphi$ & Stereographic $S^2$, position-dependent coefficients \\
\texttt{cylindrical\_kg/} & 3+1D & $\varphi$ & Mixed curved $(r,\theta)$ and flat $(z)$ directions \\
\texttt{cylindrical\_kg\_1d/} & 1+1D & $\varphi$ & Cylindrical coordinates, plane-wave 1D reduction \\
\texttt{spherical\_kg\_1d/} & 1+1D & $\varphi$ & Spherical coordinates, plane-wave 1D reduction \\
\midrule
\multicolumn{4}{@{}l}{\emph{Linearized gravity and tensor fields}} \\
\texttt{gravitational\_waves/} & 3+1D & $h_{\mu\nu}$ & xPert~\cite{brizuela2009xpert} linearization, TT gauge, rank-2 tensor \\
\texttt{gravitational\_waves\_1d/} & 1+1D & $h_{\mu\nu}$ & TT gauge + plane-wave reduction \\
\texttt{massive\_gravity/} & 2+1D & $h_{\mu\nu}$ & Fierz--Pauli mass, xPert~\cite{brizuela2009xpert}, 5 constraints \\
\texttt{massive\_3form/} & 3+1D & $C_{\mu\nu\rho}$ & Rank-3 antisymmetric, symmetry reduction $64 \to 4$ \\
\midrule
\multicolumn{4}{@{}l}{\emph{Multi-field perturbation and Gertsenshtein}} \\
\texttt{gertsenshtein/} & 1+1D & $h_\times$, $a_y$ & Einstein--Maxwell graviton-photon conversion \\
\texttt{torsion\_gertsenshtein/} & 3+1D & $h_{\mu\nu}$, $t_{\alpha\mu\nu}$, $a_\mu$ & Combined PGT+EM, torsion-independence proven \\
\texttt{graviton\_torsion/} & 3+1D & $h_{\mu\nu}$, $t_{\alpha\mu\nu}$ & Full quadratic PGT ($\alpha_I T^2 + b_5 \tilde{R}^2$), 37 fields \\
\texttt{test\_ibp/} & 1+1D & $\varphi$ & Box Lagrangian, validates IBP in canonical pipeline \\
\bottomrule
\end{tabular}
\end{table}

\subsection{Supported Operators}
\label{sec:supported-operators}

Operators map JSON terms to numpy spatial operations (2nd-order finite-difference stencils). All support cross-field references and periodic/Dirichlet/Neumann boundary conditions.

\begin{table}[htbp]
\centering
\caption{Supported spatial operators.}
\label{tab:operators}
\begin{tabular}{@{}lll@{}}
\toprule
Operator & Min Dim & Description \\
\midrule
\texttt{identity} & 1D & Field itself (mass terms) \\
\texttt{laplacian} & 1D & Full spatial Laplacian \\
\texttt{laplacian\_x}, \texttt{\_y}, \texttt{\_z} & 1D/2D/3D & Directional second derivative \\
\texttt{gradient\_x}, \texttt{\_y}, \texttt{\_z} & 1D/2D/3D & Directional first derivative \\
\texttt{cross\_derivative\_xy}, \texttt{\_xz}, \texttt{\_yz} & 2D/3D & Mixed partial derivative \\
\texttt{first\_derivative\_t} & 1D & Time derivative (friction/damping terms) \\
\texttt{biharmonic} & 1D & Fourth-order Laplacian \\
\texttt{mixed\_T\_S1\_S2\_\ldots} & 1D & Mixed time-space operators (auto-extracted from Wolfram) \\
\texttt{derivative\_N\_x} & 1D & $N$th-order spatial derivative (e.g., \texttt{derivative\_3\_x} for $\partial_x^3$) \\
\bottomrule
\end{tabular}
\end{table}

\subsection{Gradient Energy Measurement}
\label{sec:gradient-energy-measurement}

\texttt{compute\_system\_energy} in \texttt{tidal/measurement/\_energy.py} evaluates the spatially-averaged Hamiltonian from each theory's \texttt{hamiltonian\_terms}. For gradient-product terms $\langle \partial_a f \cdot \partial_b g \rangle$ the code auto-dispatches between three paths:

\begin{table}[htbp]
\centering
\caption{Gradient-energy paths in \texttt{\_evaluate\_single\_hamiltonian\_term}. The dispatch keys only on data properties (periodicity, coefficient structure); it does \emph{not} consult any solver-side flag. Measurement accuracy is decoupled from the solver mode used to produce the fields.}
\label{tab:gradient-energy-paths}
\begin{tabular}{@{}p{0.20\textwidth}p{0.35\textwidth}p{0.38\textwidth}@{}}
\toprule
Path & Trigger & Error \\
\midrule
Parseval (k-space) & all axes periodic + scalar coefficient + scalar volume weight & Machine precision (exact for band-limited periodic fields) \\
FD IBP & periodic axis with FD solver, position-dependent coefficient, or non-scalar volume weight & $O((k \cdot dx)^2 / 12)$ for 2nd-order stencil; scales as $1/N^2$ \\
Spectral-direct & spectral solver mode (legacy, retained for non-Parseval paths) & $O(1/N)$ at Nyquist due to rfft phase handling \\
\bottomrule
\end{tabular}
\end{table}

The Parseval path computes
\[
\langle \partial_a f \cdot \partial_b g \rangle \;=\; \frac{1}{\prod_i N_i} \sum_{\vec{k}} w_{\vec{k}}\, k_a k_b\, \operatorname{Re}(\hat{f}^*_{\vec{k}}\, \hat{g}_{\vec{k}})
\]
over \texttt{rfftn}'s half-spectrum, with $w_{\vec{k}} = 2$ on interior last-axis bins and $w_{\vec{k}} = 1$ at DC and (for even $N$) Nyquist. When the Parseval path is \emph{not} applicable on an all-periodic domain --- typically because a coefficient is position-dependent --- a one-shot \texttt{RuntimeWarning} signals that the term is being evaluated at FD accuracy rather than machine precision. See issue \#312 for the design rationale and the follow-up tickets for the modal-solver Fourier-coefficient shortcut and the legacy \texttt{\_use\_spectral} flag refactor.

\subsection{Component-Level Euler--Lagrange}
\label{sec:component-el}

Since v0.15.0, the default derivation strategy for all theories is \textbf{component-level Euler--Lagrange}, which operates in two phases:

\paragraph{Phase~A: Component decomposition.}
The abstract Lagrangian $\lag(g_{\mu\nu}, \phi^A, \nabla_\mu \phi^A, \ldots)$ is decomposed to a scalar function of coordinate components:
\begin{equation}
  \lag_{\text{comp}} = \lag(g_{00}, g_{01}, \ldots, \phi_0, \phi_1, \ldots, \partial_t \phi_0, \partial_x \phi_0, \ldots)
\end{equation}
This reuses the existing xCoba basis projection (\texttt{ToBasis}), metric contraction, and Christoffel evaluation stages of the Wolfram pipeline. For a 2+1D theory, this typically completes in $\sim$4\,s.

\paragraph{Phase~B: Standard calculus E--L.}
The equations of motion are derived by direct differentiation of $\lag_{\text{comp}}$ using the multi-index Euler--Lagrange formula:
\begin{equation}
  \frac{\delta \lag}{\delta q_i} = \sum_{\alpha} (-1)^{|\alpha|}\, D^\alpha \frac{\partial \lag}{\partial (D^\alpha q_i)} = 0
  \label{eq:component-el}
\end{equation}
where $D^\alpha = \partial_t^{n_t} \partial_x^{n_x} \cdots$ ranges over all derivative orders present in $\lag_{\text{comp}}$. This produces one equation per component field in under 1\,s per field.

\paragraph{Motivation.}
The previous approach---abstract \texttt{VarD} followed by \texttt{TraceBasisDummy} for index contraction---suffers from combinatorial explosion when the Lagrangian contains many contracted abstract indices. For example, the $\tilde{R}^2$ curvature-squared torsion Lagrangian has $\sim$45 contracted abstract indices, leading to $O(\text{dim}^{n})$ scaling in \texttt{TraceBasisDummy} ($>$77\,min for 2+1D, infeasible for 3+1D). Component-level E--L bypasses the abstract index machinery entirely, reducing the same derivation to $\sim$5\,s---a \textbf{900$\times$ speedup}. It is now the default for all theories, not only those with higher-curvature terms.

\paragraph{Implementation.}
\texttt{ComponentEulerLagrange} in \texttt{ComponentDecompose.wl} handles the Phase~B computation, including integration by parts and support for antisymmetric tensor field components via symmetry-permuted tuple replacement (\texttt{ReplaceHigherRankFieldComponents}).

\subsection{Perturbative Reduction (v6)}
\label{sec:ostrogradsky}

Higher-derivative Lagrangians (e.g., $R^2$, $\tilde{R}^2$ curvature-squared terms) produce equations of motion with time derivatives of order $>2$. Since all PDE solver backends support at most second-order-in-time equations, TIDAL applies the iterative \textbf{perturbative reduction} pipeline introduced in v6 to evolve only the physical (ghost-free) solution branch.

The complete mathematical and implementation treatment lives in \cref{sec:perturbative-reduction} (\texttt{perturbative\_reduction.tex}). A summary follows for architectural context.

\paragraph{Pipeline.} Theories that contain higher-derivative terms declare a \verb|[perturbation]| section in their \verb|theory.toml| naming the small parameters (e.g. \verb|small_parameters = ["b5"]|). At derive time, each emitted \verb|OperatorTerm| is tagged with \verb|order_in_eps| = total exponent of those small parameters in its symbolic coefficient. At simulate time the \verb|PerturbativeSolver| orchestrates:
\begin{enumerate}
  \item \textbf{Pass 0.} \verb|EquationSystem.base_spec(small_parameters)| sets every small parameter to zero, demoting any equation whose kinetic coefficient then vanishes to an algebraic constraint (time order 0). The resulting ghost-free second-order system is solved by the standard modal solver with \verb|return_eigendata=True|.
  \item \textbf{Pass 1.} The order-1 correction source $M_{\mathrm{src}}(\mathbf{k}) \cdot y^{(0)}(t)$ is evolved by a closed-form Duhamel integral that reuses the Pass 0 eigendecomposition. Constraint fields are recovered at both orders via the existing Schur operator $\mathsf{S} = -\mathsf{S}_{cc}^{-1} \mathsf{S}_{cd}$.
  \item \textbf{Combined.} $q_{\mathrm{total}} = \sum_n q^{(n)}$ returned alongside a validity diagnostic $\max(\varepsilon \cdot |\lambda|^2 \cdot t_{\mathrm{end}})$ that flags approaching EFT breakdown.
\end{enumerate}

\paragraph{CLI.} The flag \verb|--perturbative-order N| defaults to~$1$ when the JSON carries a \verb|perturbation| metadata block and~$0$ otherwise. Theories without the metadata block simulate through the plain modal path unchanged. The flag is silent about orders beyond~$1$; see the scope limits in \cref{sec:pr-v6}.

\paragraph{Quartic EM and matter-only derivative-only theories.} Euler--Heisenberg $(F_{\mu\nu}F^{\mu\nu})^2$ and any single-dynamical-field theory whose field enters only through covariant derivatives (e.g.\ a photon via $F = \mathrm{d}A$) are supported as of v0.33.9; see \cref{sec:pr-eh-cd} for the upstream (pre-xPert) Lagrangian-assignment fix (Hold + Scalar-wrapping against Power auto-distribution) and the correctness-aware CD ComponentValue precompute gate that unblocks these theories.

\paragraph{Ghost-free by construction.} The Pass 0 evolution operator is always the second-order base; the higher-derivative terms appear only as sources on the right-hand side. No ghost modes propagate because they are not part of the solution ansatz. This is the central architectural guarantee of the v6 path.

\paragraph{Historical note.} Before v6, TIDAL used a mechanical Ostrogradsky reduction that introduced auxiliary fields $w_q = \partial_t^2 q$ to flatten higher-order equations into coupled second-order ones. The reduction was theoretically sound but produced ghost modes whose long-time growth corrupted simulations. It was replaced by the iterative path in v6 (\cref{sec:pr-v6}) and the \verb|tidal/symbolic/ostrogradsky.py| module was removed.
